\section{结果与展示材料}
\begin{strip}
\centering
\captionof{table}{真机任务成功率（\%）。OOD 汇总未见桌布、物体外观和干扰物三类条件；平均值对四个任务等权计算。}
\label{tab:real-results}
\small
\setlength{\tabcolsep}{5pt}
\renewcommand{\arraystretch}{1.12}
\begin{tabularx}{\textwidth}{@{}l*{10}{>{\centering\arraybackslash}X}@{}}
\toprule
\multirow{2}{*}{方法} & \multicolumn{2}{c}{放入抽屉} & \multicolumn{2}{c}{放进盒子} & \multicolumn{2}{c}{积木分类} & \multicolumn{2}{c}{插接} & \multicolumn{2}{c}{平均} \\
\cmidrule(lr){2-3}\cmidrule(lr){4-5}\cmidrule(lr){6-7}\cmidrule(lr){8-9}\cmidrule(l){10-11}
 & ID & OOD & ID & OOD & ID & OOD & ID & OOD & ID & OOD \\
\midrule
OpenWAM-Alpha & \blankresult & \blankresult & \blankresult & \blankresult & \blankresult & \blankresult & \blankresult & \blankresult & \blankresult & \blankresult \\
X-WAM & \blankresult & \blankresult & \blankresult & \blankresult & \blankresult & \blankresult & \blankresult & \blankresult & \blankresult & \blankresult \\
\rowcolor{gray!10}
\textbf{MetisWAM4D} & \blankresult & \blankresult & \blankresult & \blankresult & \blankresult & \blankresult & \blankresult & \blankresult & \blankresult & \blankresult \\
\bottomrule
\end{tabularx}
\normalsize

\end{strip}

\subsection{正文论述的填空顺序}
\paragraph{总体表现。}填写 ID 及三类 OOD 的各任务成功率，随后报告任务平均值与相对两个基线的百分点差。每项平均值沿用同一组任务和同一权重。

\paragraph{外观变化。}分别比较桌布、物体外观、干扰物三列，并记录相对于 ID 的变化。结论围绕实际改善的条件展开，附代表性成功和失败片段。

\paragraph{精细操作。}对照普通放置与插接、抽屉边界操作的提升量。补充主要失败原因，如抓取偏差、边界碰撞、对齐失败、目标混淆和超时。

\subsection{结果记录字段}
每次试验保存任务、方法、checkpoint、条件、初始化编号、开始／结束时间、成功标记、结束原因与录像路径。成功率由逐次记录汇总；表格保留成功数和总次数，完成时间可在附录单独补充。

\begin{figure}[htbp]
\centering
\panel{28mm}{OOD 场景占位\\ID \quad 桌布 \quad 物体外观 \quad 干扰物\\同一任务、同一相机、相近初始位置}
\caption{条件对照图占位。每列标注实际改变的因素与道具编号。}
\end{figure}

\subsection{机制与视频材料}
实机执行录像用于展示策略完成任务的过程。若展示模型内部预测，再补充预测 RGB、预测轨迹与实际执行的对应片段；分支训练方式、预测步长和采样配置随图记录。

\begin{figure}[htbp]
\centering
\panel{31mm}{可选机制图占位\\当前 RGB／深度／mask\\预测 RGB／Track\\真实执行关键帧}
\caption{几何和运动表示展示占位。参考 FlowWAM Figure 9 的预测与执行并排结构。}
\end{figure}
